% TOP_PROBLEM: {"id": 9700040, "title": "Stationary law of a drift-jump particle process", "category_id": 19, "category": {"id": 19, "name": "probability", "display_name": "Probability", "description": "Problems involving probability theory, stochastic processes, and random structures.", "slug": "probability", "order_index": 19, "created_at": "2026-07-31T15:26:25.671Z"}, "status": "open", "problem_number": "AMR-096-0040", "set_id": 13}
% FIRST_POSTED: 2026-10-01
% ENTRY_KIND: statement_only
% UnsolvedMath ID 9700040; source catalogue code AMR-096-0040
% !TeX program = lualatex
% Definition and sources only; this file is not an LLM solution attempt.
\documentclass[11pt,a4paper]{article}
\usepackage[margin=25mm]{geometry}
\usepackage{fontspec}
\setmainfont{lmroman10-regular.otf}[BoldFont=lmroman10-bold.otf,
 ItalicFont=lmroman10-italic.otf,BoldItalicFont=lmroman10-bolditalic.otf]
\usepackage{amsmath,amssymb,mathtools}
\usepackage{unicode-math}
\setmathfont{latinmodern-math.otf}
\AtBeginDocument{\let\setminus\smallsetminus}
\usepackage{xurl,hyperref}
\hypersetup{colorlinks=true,urlcolor=blue}
\setcounter{secnumdepth}{0}
\setlength{\parindent}{0pt}
\setlength{\parskip}{5pt}
\setlength{\emergencystretch}{3em}
\begin{document}
\section*{Stationary law of a drift-jump particle process}\label{9700040}
{\small UnsolvedMath ID 9700040 · AMR-096-0040 · Probability · AMR Open Problem Lists}
\subsection{Definitions and mathematical statement}
Fix $k\ge1$ and consider ordered particles $0<X_1(t)<\cdots<X_k(t)$ on the positive half-line. Between jumps each follows $dX_i/dt=X_i$, so a position grows exponentially during an interval without jumps. Independently, sample a Poisson point process on $(0,\infty)\times(0,\infty)$ with intensity $dx\,dt$.

At an event $(u,t)$, move the leftmost particle strictly to the right of $u$ instantaneously to $u$; if none exists, do nothing. The ordering of the particles is preserved. More explicitly, with $x_0=0$, the generator on differentiable test functions is
\[
\mathcal L f(x)=\sum_{i=1}^k x_i\partial_i f(x)
+\sum_{i=1}^k\int_{x_{i-1}}^{x_i}
[f(x_1,\ldots,x_{i-1},u,x_{i+1},\ldots,x_k)-f(x)]\,du.
\]
The first $j$ particles evolve autonomously with the same rule for every $j<k$, giving consistent finite-particle descriptions of an infinite system.

Find a reasonably explicit description of its stationary law. In concrete terms, determine the invariant joint distributions of $(X_1,\ldots,X_k)$ for every $k$, consistently in $k$, for example through densities, transforms, or a direct sampling representation.

Stationarity means that starting with this distribution leaves it unchanged at every deterministic time. The source supplies uniqueness as background; the task is to identify the distribution, not merely prove that some invariant measure exists. The drift and the jump intensity above fix the normalization of both space and time.
\subsection{Short English statement}
Can the stationary joint distributions of this exponentially drifting, left-jumping particle system be described explicitly?
\subsection{Sources}
\begin{itemize}
\item[S1] ULAM.AI and the UnsolvedMath Contributors, supplied problems.json, record 9700040 (AMR-096-0040), AMR Open Problem Lists. Catalogue identity and source transcription; CC BY 4.0.
\url{https://huggingface.co/datasets/ulamai/UnsolvedMath}.
\item[S2] Aldous - Open problems, left\_Hamm.html. Original source indexed by AMR-096-0040.
\url{https://www.stat.berkeley.edu/~aldous/Research/OP/index.html}.
\item[S3] D. Aldous, A one-dimensional drift-jump particle process, precise finite-particle dynamics.
\url{https://www.stat.berkeley.edu/~aldous/Research/OP/left_Hamm.html}.
\end{itemize}
\end{document}
